% 第 6 章：三角函数
% 说明：本章按内容编排，课时划分由教师把握。
\chapter{三角函数}

\begin{objective}
\begin{itemize}
  \item 理解任意角的概念，会进行角度与弧度的互化，会计算弧长与扇形面积；
  \item 会用单位圆定义任意角的正弦、余弦、正切，会求特殊角的三角函数值；
  \item 掌握同角三角函数关系与诱导公式，会化简与求值；
  \item 会画正弦、余弦、正切的图象，说出它们的周期性、奇偶性与单调性；
  \item 理解 $y=A\sin(\omega x+\varphi)$ 中 $A,\omega,\varphi$ 的意义，会作图象变换；
  \item 会用三角函数建立简单的周期现象模型。
\end{itemize}
\end{objective}

\begin{keypoint}
重点：任意角与弧度制；单位圆定义；同角关系与诱导公式；正弦、余弦、正切的图象与性质；
$y=A\sin(\omega x+\varphi)$ 的图象变换。\\
难点：弧度制的意义；诱导公式的“符号看象限”；用五点法作图；振幅、周期、初相的综合。\\
课时建议：本章建议 $5\sim6$ 课时（角与弧度 $1$、三角函数定义与诱导公式 $1\sim2$、
图象与性质 $1\sim2$、$y=A\sin(\omega x+\varphi)$ 与模型应用 $1$），由教师根据学情把握。
\end{keypoint}

\begin{procedure}

\section{任意角与弧度制}
\textit{温故}：初中锐角三角函数（直角三角形中的边的比）；特殊角 $30^\circ$、$45^\circ$、$60^\circ$ 的正弦值。

\textit{新知}：角看成射线绕端点旋转形成的图形，逆时针为正、顺时针为负。
终边相同的角：$\{\beta\mid\beta=\alpha+k\cdot360^\circ,\ k\in\mathbb{Z}\}$。
判断象限：先化到 $0^\circ\sim360^\circ$。

弧度制定义：长度等于半径的弧所对的圆心角为 $1\ \text{rad}$。
换算：$180^\circ=\pi\ \text{rad}$，$1^\circ=\frac{\pi}{180}\ \text{rad}$。
弧长 $l=|\alpha|r$，扇形面积 $S=\frac12lr=\frac12|\alpha|r^2$（$\alpha$ 用弧度）。

\key{为什么用弧度}：与半径无关，公式更简洁，且使三角函数的导数公式干净；
高等数学与物理一律用弧度。

\begin{example}
把 $60^\circ$、$-135^\circ$ 化成弧度，把 $\dfrac{2\pi}{3}$ 化成角度；
已知扇形半径为 $2$、圆心角 $\dfrac{\pi}{3}$，求弧长与面积。
\end{example}
\begin{solution}
$60^\circ=\dfrac{\pi}{3}$，$-135^\circ=-\dfrac{3\pi}{4}$，$\dfrac{2\pi}{3}=120^\circ$；
$l=\dfrac{\pi}{3}\times2=\dfrac{2\pi}{3}$，$S=\dfrac12lr=\dfrac{2\pi}{3}$。
\end{solution}

\section{任意角的三角函数}
单位圆定义：设角 $\alpha$ 的终边与单位圆交于 $P(x,y)$，则
$\sin\alpha=y$，$\cos\alpha=x$，$\tan\alpha=\dfrac{y}{x}$（$x\neq0$）。
一般地，终边过 $P(x,y)$，$r=\sqrt{x^2+y^2}$，则 $\sin\alpha=\frac{y}{r}$，$\cos\alpha=\frac{x}{r}$，$\tan\alpha=\frac{y}{x}$。

同角关系：$\sin^2\alpha+\cos^2\alpha=1$，$\tan\alpha=\dfrac{\sin\alpha}{\cos\alpha}$。
诱导公式：$\sin(-\alpha)=-\sin\alpha$，$\cos(-\alpha)=\cos\alpha$，
$\sin(\pi-\alpha)=\sin\alpha$，$\cos(\pi-\alpha)=-\cos\alpha$，
$\sin(\pi+\alpha)=-\sin\alpha$，$\cos(\pi+\alpha)=-\cos\alpha$，
$\sin\left(\frac{\pi}{2}-\alpha\right)=\cos\alpha$，$\cos\left(\frac{\pi}{2}-\alpha\right)=\sin\alpha$。
口诀：\key{奇变偶不变，符号看象限}。

\begin{example}
已知 $\sin\alpha=\dfrac35$，$\alpha$ 为第二象限角，求 $\cos\alpha$ 与 $\tan\alpha$。
\end{example}
\begin{solution}
$\cos^2\alpha=1-\frac{9}{25}=\frac{16}{25}$，第二象限 $\cos\alpha<0$，故 $\cos\alpha=-\frac45$；
$\tan\alpha=\frac{\sin\alpha}{\cos\alpha}=-\frac34$。
\end{solution}

\begin{exercise}
\begin{enumerate}
  \item 已知角 $\alpha$ 终边过点 $(-4,3)$，求 $\sin\alpha$、$\cos\alpha$、$\tan\alpha$；
  \item 化简 $\sin(180^\circ+\alpha)\cos(-\alpha)$，并求 $\sin\dfrac{7\pi}{6}$。
\end{enumerate}
\end{exercise}

\section{三角函数的图象与性质}
\textit{五点法}：取 $x=0$、$\dfrac{\pi}{2}$、$\pi$、$\dfrac{3\pi}{2}$、$2\pi$ 五个点。
\begin{itemize}
  \item $y=\sin x$：定义域 $\mathbb{R}$，值域 $[-1,1]$，周期 $2\pi$，奇函数；
        在 $\left[-\dfrac{\pi}{2}+2k\pi,\dfrac{\pi}{2}+2k\pi\right]$ 上递增。
  \item $y=\cos x$：定义域 $\mathbb{R}$，值域 $[-1,1]$，周期 $2\pi$，偶函数；
        在 $[2k\pi,\pi+2k\pi]$ 上递减。
  \item $y=\tan x$：定义域 $\left\{x\mid x\neq\dfrac{\pi}{2}+k\pi\right\}$，值域 $\mathbb{R}$，
        周期 $\pi$，奇函数。
  \item 它在每个 $\left(-\dfrac{\pi}{2}+k\pi,\dfrac{\pi}{2}+k\pi\right)$ 内递增，
        渐近线 $x=\dfrac{\pi}{2}+k\pi$。
\end{itemize}

\begin{figure}[htbp]
  \centering
  \begin{tikzpicture}
    \begin{axis}[
      width=0.80\textwidth, height=6.4cm,
      xmin=-2*pi, xmax=2*pi, ymin=-1.6, ymax=1.6,
      xtick={-2*pi, -3*pi/2, -pi, -pi/2, 0, pi/2, pi, 3*pi/2, 2*pi},
      xticklabels={$-2\pi$,$-\frac{3\pi}{2}$,$-\pi$,$-\frac{\pi}{2}$,$0$,$\frac{\pi}{2}$,$\pi$,$\frac{3\pi}{2}$,$2\pi$},
      trig format plots=rad,
      xlabel=$x$, ylabel=$y$,
      grid=major, grid style={dashed,gray!35},
      legend pos=north east, legend style={font=\footnotesize},
      tick label style={font=\footnotesize},
      title={\small $y=\sin x$ 与 $y=\cos x$（周期 $2\pi$）},
      title style={font=\small},
    ]
      \addplot[domain=-2*pi:2*pi, samples=240, thick, main!80!black] {sin(x)};
        \addlegendentry{$y=\sin x$}
      \addplot[domain=-2*pi:2*pi, samples=240, thick, second] {cos(x)};
        \addlegendentry{$y=\cos x$}
      \addplot[only marks, mark=*, mark size=1.4pt, black]
        coordinates {(0,0) (pi/2,1) (pi,0) (3*pi/2,-1) (2*pi,0)};
    \end{axis}
  \end{tikzpicture}
  \caption{正弦与余弦：黑点是 $y=\sin x$ 的五点法关键点}
\end{figure}

\begin{figure}[htbp]
  \centering
  \begin{tikzpicture}
    \begin{axis}[
      width=0.80\textwidth, height=6.4cm,
      xmin=-4.8, xmax=4.8, ymin=-4.5, ymax=4.5,
      xtick={-3*pi/2, -pi, -pi/2, 0, pi/2, pi, 3*pi/2},
      xticklabels={$-\frac{3\pi}{2}$,$-\pi$,$-\frac{\pi}{2}$,$0$,$\frac{\pi}{2}$,$\pi$,$\frac{3\pi}{2}$},
      trig format plots=rad,
      xlabel=$x$, ylabel=$y$,
      grid=major, grid style={dashed,gray!35},
      legend pos=north west, legend style={font=\footnotesize},
      tick label style={font=\footnotesize},
      title={\small $y=\tan x$：周期 $\pi$，虚线为渐近线},
      title style={font=\small},
    ]
      \addplot[domain=-1.32:1.32, samples=120, thick, main!80!black] {tan(x)};
      \addplot[domain=1.83:4.45, samples=120, thick, main!80!black] {tan(x)};
      \addplot[domain=-4.45:-1.83, samples=120, thick, main!80!black] {tan(x)};
      \addplot[dashed, gray] coordinates {(-pi/2,-4.5) (-pi/2,4.5)};
      \addplot[dashed, gray] coordinates {(pi/2,-4.5) (pi/2,4.5)};
      \addplot[dashed, gray] coordinates {(-3*pi/2,-4.5) (-3*pi/2,4.5)};
      \addplot[dashed, gray] coordinates {(3*pi/2,-4.5) (3*pi/2,4.5)};
      \addlegendentry{$y=\tan x$}
    \end{axis}
  \end{tikzpicture}
  \caption{正切曲线：每个周期内递增，渐近线处断开}
\end{figure}

\section{函数 $y=A\sin(\omega x+\varphi)$ 的图象与性质}
三个参数：$|A|$ 为振幅（纵向）；周期 $T=\dfrac{2\pi}{\omega}$（横向）；初相 $\varphi$（左右平移 $\frac{\varphi}{\omega}$）。
变换顺序：平移 $\to$ 横缩 $\to$ 纵伸；或代入五个关键点。

\begin{example}
求 $y=2\sin\left(2x+\dfrac{\pi}{3}\right)$ 的振幅、周期、初相，并说明变换过程。
\end{example}
\begin{solution}
$A=2$，$T=\pi$，$\varphi=\frac{\pi}{3}$。
$y=\sin x$ 左移 $\frac{\pi}{3}$ 得 $y=\sin\left(x+\frac{\pi}{3}\right)$；
横坐标缩短为原来的 $\frac12$ 得 $y=\sin\left(2x+\frac{\pi}{3}\right)$；
纵坐标伸长为原来的 $2$ 倍。
\end{solution}

\begin{figure}[htbp]
  \centering
  \begin{tikzpicture}
    \begin{axis}[
      width=0.80\textwidth, height=6.4cm,
      xmin=0, xmax=2*pi, ymin=-2.6, ymax=2.6,
      xtick={0, pi/2, pi, 3*pi/2, 2*pi},
      xticklabels={$0$,$\frac{\pi}{2}$,$\pi$,$\frac{3\pi}{2}$,$2\pi$},
      trig format plots=rad,
      xlabel=$x$, ylabel=$y$,
      grid=major, grid style={dashed,gray!35},
      legend pos=north east, legend style={font=\footnotesize},
      tick label style={font=\footnotesize},
      title={\small $y=\sin x$ 与 $y=2\sin\left(2x+\frac{\pi}{3}\right)$},
      title style={font=\small},
    ]
      \addplot[domain=0:2*pi, samples=200, thick, gray] {sin(x)};
        \addlegendentry{$y=\sin x$}
      \addplot[domain=0:2*pi, samples=400, thick, main!80!black] {2*sin(2*x+pi/3)};
        \addlegendentry{$y=2\sin\left(2x+\frac{\pi}{3}\right)$}
    \end{axis}
  \end{tikzpicture}
  \caption{振幅变大、周期变小、整体左移：同一族曲线形状相同，位置与大小不同}
\end{figure}

\begin{exercise}
求 $y=3\sin\left(\dfrac{x}{2}-\dfrac{\pi}{4}\right)$ 的振幅、周期、初相，并说明变换过程。
\end{exercise}

\section{三角函数模型的简单应用}
周期现象可用 $y=A\sin(\omega x+\varphi)+b$ 刻画，$b$ 是平衡位置，$A$ 是最大偏离。
例：温度 $T(t)=8\sin\left(\frac{\pi}{12}t-\frac{7\pi}{12}\right)+20$，求最高、最低温度与周期。

\key{数学文化}：三角学起源于测量（古希腊弦表、印度正弦、阿拉伯系统化），
17 世纪后成为分析学的一部分；傅里叶级数说明任何周期现象都可用三角函数叠加逼近。

\begin{exercise}
弹簧振子 $y=5\sin(2\pi t)$（cm，s），求振幅、周期与 $t=0.25$ s 时的位移。
\end{exercise}

\end{procedure}

\begin{board}
\begin{itemize}
  \item 角度与弧度：$180^\circ=\pi$ rad；弧长 $l=|\alpha|r$，面积 $S=\frac12|\alpha|r^2$；
  \item 三角函数定义（单位圆）：$\sin\alpha=y$，$\cos\alpha=x$，$\tan\alpha=\frac{y}{x}$；
  \item 同角关系：$\sin^2\alpha+\cos^2\alpha=1$，$\tan\alpha=\frac{\sin\alpha}{\cos\alpha}$；诱导公式口诀“奇变偶不变，符号看象限”；
  \item $y=\sin x$、$y=\cos x$ 周期 $2\pi$；$y=\tan x$ 周期 $\pi$；
  \item $y=A\sin(\omega x+\varphi)$：振幅 $|A|$、周期 $\frac{2\pi}{\omega}$、初相 $\varphi$；
  \item 周期现象模型：$y=A\sin(\omega x+\varphi)+b$。
\end{itemize}
\end{board}

\begin{homework}
\begin{enumerate}
  \item 把 $150^\circ$、$-210^\circ$ 化成弧度；把 $\dfrac{5\pi}{6}$ 化成角度；
  \item 写出与 $-60^\circ$ 终边相同的角的集合，并找出 $-360^\circ\sim0^\circ$ 内的那个角；
  \item 已知 $\cos\alpha=-\dfrac{5}{13}$，$\alpha$ 在第三象限，求 $\sin\alpha$、$\tan\alpha$；
  \item 化简 $\dfrac{\sin(2\pi-\alpha)\cos(\pi+\alpha)}{\tan(-\alpha)}$；
  \item 用五点法在 $[0,2\pi]$ 上画 $y=1+2\sin x$ 的图象（写出五个关键点）；
  \item 判断 $y=\sin x+\cos x$ 的奇偶性，并求最大值；
  \item 函数 $y=A\sin(\omega x+\varphi)$ 相邻两条最高点距离为 $\pi$，过点 $\left(0,\frac12\right)$，振幅为 $1$，写出一个解析式；
  \item 阅读“数学文化”，举例说明三角函数在生活中的一个周期现象。
\end{enumerate}
\end{homework}
